\section{Introduction}
\subsection{Scope}

This document describes how to reduce data with ESO Data Processing System (EDPS) from the command line. It is meant for first time users of EDPS, and guides through the installation procedure and data reduction using the ESPRESSO pipeline and its demo data as an example. After working through this tutorial, the reader should be able to use other supported EDPS workflows for the reduction of user provided data.

A graphical interface (EDPS-GUI) is also available; installation and usage are described at \url{https://www.eso.org/sci/software/edps.html}. Developers who need to write or modify workflows should consult the EDPS workflow design tutorial.

\subsection{What is EDPS?}

The ESO Data Processing System (EDPS) is a framework  to  run ESO's data processing pipelines and it is meant to eventually replace the previous EsoReflex environment (\url{https://www.eso.org/sci/software/esoreflex/}). The general principles of EDPS have been described by Freudling, Zampieri, Coccato et al. \href{https://www.aanda.org/articles/aa/full_html/2024/01/aa47651-23/aa47651-23.html}{(2024, A\&A, 681, A93)}. Please refer to that paper if you have used EDPS for research resulting in a scientific publication. 

Each of ESO's data processing pipeline consist of a series of standalone programs called "recipes". Each recipe is designed to process certain type(s) of input data. The processing 
of these input data typically requires a range of auxiliary files such as calibration files. EDPS 
is designed to select appropriate input data for the different recipes of a pipeline, and execute 
them in sequence. This is done by specifying for each pipeline the workflow for organising data
and executing the recipes. This workflow can be used to process a set of data fully automatically. 
%either in batch or interactive mode. 


